\documentclass[12pt,a4paper,onecolumn]{article}
\title{Digest of The New York Times}
\author{The New York Times}
% \definecolor{dark-green}{rgb}{0,0.25,0}
% \usepackage[utf8]{inputenc}

\usepackage{fontspec,xunicode,xltxtra}
\setmainfont[Ligatures=Required,Mapping=tex-text]{Georgia}
% \setmainfont[Ligatures=Required,Mapping=tex-text]{Palatino Linotype}
\setsansfont[Ligatures=Required,Mapping=tex-text]{Myriad Pro}
\setmonofont[Ligatures=Required,Mapping=tex-text]{Courier New}
\usepackage{xeCJK}
\setCJKmainfont{宋体}
\setCJKmonofont{黑体}

% \usepackage{indentfirst}


\usepackage[dvips,dvipsnames,svgnames]{xcolor}
\definecolor{light-gray}{gray}{0.95}

%生成PDF的链接
\usepackage[dvipdfmx]{hyperref}
\hypersetup{
    bookmarks=true,         % show bookmarks bar?
    bookmarksopen=true,
    pdfpagemode=UseNone,    % options: UseNode, UseThumbs, UseOutlines, FullScreen
    pdfstartview=FitH,
    pdfborder=1,
    pdfhighlight=/P,
    pdfauthor={wuxch},
    unicode=true,           % non-Latin characters in Acrobat’s bookmarks
    colorlinks,             % false: boxed links; true: colored links
    linkcolor=blue,         % color of internal links
    citecolor=blue,        % color of links to bibliography
    filecolor=magenta,      % color of file links
    urlcolor=cyan           % color of external links
}
% 使用目录
\usepackage{makeidx}
% 使用图形
\usepackage{graphicx}


\usepackage{titletoc}
\dottedcontents{section}[1.5em]{}{2.3em}{1pc}

% % 页眉页脚的设定
% \usepackage{fancyhdr}
% \pagestyle{fancy}
% \lhead{\includegraphics[height=0.2in]{zte_logo}}

% \usepackage{footmisc}
\usepackage{pfnote}
\usepackage{fnpos}

% % 嵌入的代码显示
% \usepackage{listings}
% \lstset{language=C++, breaklines, extendedchars=false}
% \lstset{keywordstyle=\color{blue},
%         commentstyle=\color{SeaGreen},
%         stringstyle=\ttfamily,
%         showstringspaces=false,
%         backgroundcolor=\color{light-gray}}
% \lstset{basicstyle=\small}

\usepackage{booktabs}

\usepackage{tipa}
\usepackage{phonetic}
\usepackage[gen]{eurosym}

% \usepackage{CJK}
\usepackage{makeidx}
% \usepackage{glossary}

\usepackage{fancyhdr}

\usepackage{lastpage}
\cfoot{\textsf{\thepage\ of \pageref*{LastPage}}}

\pagestyle{fancy}
\fancyhead{}
\renewcommand{\headrulewidth}{0pt}
\pagenumbering{arabic}

\usepackage{graphicx}
\usepackage{picinpar}

\setlength{\parindent}{0in}
\setlength{\parskip}{2.5ex plus4mm minus3mm}
% \setlength{\parskip}{.3\baselineskip}

\setcounter{section}{0}

% bibtex
\usepackage[super]{natbib}
% \usepackage[globalcitecopy,labelstoglobalaux]{bibunits}
% \makeatletter
% \renewcommand{\@bibunitname}{\jobname.\the\@bibunitauxcnt}
% \makeatother

% \renewcommand{\refname}{缩略语}
% \renewcommand{\refname}{}

\makeatletter                   % Reference list option change
\renewcommand\@biblabel[1]{#1:  } % from [1] to 1: 
\makeatother                    %

\newcommand{\citenumfont}[1]{\textsf{#1}}

% \usepackage{tipa}
% \usepackage{phonetic}

\newcommand{\phonetic}[1]{{\fontspec[Colour=CC00CC]{Microsoft Sans Serif} #1}}


%%% Local Variables:
%%% TeX-master: "medical_trivia"
%%% End:

\makeindex
\begin{document}
\begin{CJK*}{GBK}{song}
\date{}
% \maketitle
\tableofcontents
\pagebreak
\section{Apple Tells Shareholders That Jobs Will Return}

CUPERTINO, Calif. -- For the last nine months, Apple has refused to get into specifics about the
well-being of its chief executive, Steven P. Jobs, even as he said last month that he was taking a
six-month leave of absence to deal with health problems.

On Tuesday, the company's shareholders had their chance to press for more information -- but they
did not get far.

At its annual shareholder meeting here on Apple's corporate campus, run by the chief operating
officer, Timothy D. Cook, the company responded to inquiries about Mr.~Jobs by saying that he still
planned to return to the company in June.

``He is deeply involved in all strategic matters and has delegated day-to-day authority to Tim Cook
and his team,'' said Arthur D. Levinson, a co-lead director of Apple and the chief executive of
Genentech. ``That's where it stands.''

Mr.~Levinson said that Apple's board regularly discussed the matter of succession at the company and
that ``if there is new information we deem of import to disclose, that will happen.''

Daniel Cooperman, Apple's general counsel, declined to answer a question about a reported
investigation by securities regulators into how Apple has handled the disclosure of information
about Mr. Jobs's condition.

At the meeting, shareholders approved the re-election of Apple’s eight-member board, which includes
Mr.~Jobs, Al Gore and Google's chief executive, Eric E. Schmidt. All of the board members were
present except for Mr.~Schmidt and Mr.~Jobs.

Shareholders also voted down a proposal, offered by the American Federation of State, County and
Municipal Employees union, that would have given investors an advisory vote each year on the pay
packages for top executives. Congress is considering mandating a similar measure.

Apple's response to questions about Mr.~Jobs's condition did not satisfy some shareholders. Brandon
Rees, a representative of the A.F.L.-C.I.O. Reserve Fund, who raised the issue during a
question-and-answer session at the meeting, said afterward that the company needed to say more about
an issue that could clearly affect its stock price.

``I was disappointed by the fact that the board was not more transparent about the health of
Mr.~Jobs in addition to a succession plan,'' Mr.~Rees said.

But other shareholders said that Mr.~Jobs was entitled to some amount of privacy, even as they
lamented his absence.

``I came to see if he was here,'' said Steven Winegarden, a longtime shareholder. ``I do believe he
has made a very significant contribution over the last 12 years, and it will be very difficult for
someone else to capture the attention of everybody in the same way.''

At one point, Apple shareholders acknowledged Mr.~Jobs's 54th birthday, which was Tuesday, by
standing and singing ``Happy Birthday.''

\section{Hong Kong Court Delays PCC's Plan to Privatize}

HONG KONG -- A bid to take private the telecommunications company PCCW hit a snag on Tuesday when
Hong Kong's high court gave market regulators three weeks to investigate accusations of vote fraud
-- the latest twist in one of Hong Kong's most publicized corporate stories.

Plans by Richard Li, a prominent Hong Kong executive, to privatize PCCW have been plagued with
problems since they were announced in November. The high court's decision on Tuesday to put off its
ruling fueled speculation that the privatization could be derailed -- or at least delayed further.

The court had been expected to rubber-stamp the shareholder vote from earlier this month, at which
the privatization was approved. That vote, however, was marred by questions of possible
vote-rigging, which led Hong Kong's Securities and Futures Commission to seize voting records.

Concerns over possible manipulation emerged after David Webb, a shareholder activist and an
outspoken critic of Hong Kong's tight-knit business community, told regulators he had received an
anonymous e-mail message saying that PCCW shares had been offered to people in exchange for
supporting the buyout offer.

The Singapore-listed holding company Pacific Century Regional Developments, which is controlled by
Mr.~Li and has a stake in PCCW, said in an e-mailed statement on Tuesday that it had ``noted'' the
court's decision and reiterated that it and its directors had no knowledge of or involvement in any
improper activities.

Nevertheless, the controversy has raised concerns about corporate governance in Hong Kong,
especially as it coincides with a decision -- unrelated to PCCW -- by the Hong Kong stock exchange
earlier this month to water down strict new limits on share trading by company directors.

PCCW will have a week to respond to the Securities and Futures Commission's evidence. Another court
hearing has been set for April 1 and 2, the securities commission said in a statement. Mr.~Li's
prominence, and the fact that many ordinary Hong Kong residents bought the stock, drew wide
attention to the PCCW story .

Mr.~Li is the son of Li Ka-Shing, a billionaire who wields substantial power in business and
political circles in Hong Kong. Richard Li bought PCCW at the height of the technology bubble in
2000 for \$28.5 billion.

Since then, the value of the company has imploded, and various attempts in the last three years to
sell all or parts of PCCW have fallen through.

Last November, Mr.~Li instead announced a bid to take the company private: PCRD and the Chinese
fixed-line telecoms carrier China Netcom Group offered 4.20 Hong Kong dollars, or \$0.54, a share
and later raised the bid to 4.50 dollars for the 52.4 percent of the company they did not already
own.

The offer, totaling about \$2 billion, drew widespread criticism from investors dismayed at the
collapse in PCCW's share price and accusations that it undervalued the company.


\section{China's Answer to a Crime Includes Amateur Sleuths}

BEIJING -- As they awaited their fates in holding cell No. 9 of the Puning County jail, Li Qiaoming
and a half-dozen fellow inmates played ``elude the cat,'' a Chinese hide-and-seek that might be
better described as Marco Polo without the pool.

Mr.~Li put on a blindfold, groped for quarry and soon was lying on the floor with a fatal head
injury.

That, at least, is the story disseminated by the police in Yunnan Province, where Mr.~Li, a
24-year-old farmer, lived and died this month after his arrest on charges of illegal logging.

Public reaction to the official explanation of Mr.~Li's death was swift, and cynical. Thousands of
Internet users surmised that Mr.~Li had died from a beating by the police. Mistreatment of people in
custody is not uncommon in China. But what happened next is.

Rather than suppress the accusations by erasing the Web postings, provincial officials invited the
public to help solve the case. They sorted through 1,000 volunteers and picked a 15-member committee
that would visit the scene of the crime, cull the evidence and ``discover the truth.''

The results ultimately proved dispiriting to those who hoped for a thorough investigation of the
police. But the case riveted the public and fueled a frank discussion online and in the state-run
media about the extent -- and the limits -- of official attempts to shape popular opinion.

In explaining his motives for a citizen investigation, Wu Hao, an official in Yunnan's propaganda
department, said he was hoping to restore the public's faith in a government that could be
unresponsive and at times hostile to accusations of misconduct. In an interview with Chinese
reporters this week, he acknowledged that the authorities could have easily quashed the debate by
censoring the Internet or could have stonewalled the calls for justice.

``Past experience has shown that the doubts of the netizens will not shift or recede on their own
over time,'' he said. ``Instead, the doubts will actually rise.''

His conclusion: ``A matter of Internet public opinion must be solved by Internet methods.''

Although the reaction to Mr.~Wu's effort was initially favorable, it soured as the limits of the
inquiry became apparent. When they arrived at the jail last Friday, the committee members were given
access to the crime scene but were not allowed to view surveillance tapes, examine the autopsy
report or question the guards on duty at the time.

They were also not permitted to interview the prime suspect, Pu Huayong, an inmate who the police
said had been unhappy with the outcome of the ``elude the cat'' game. The official police report
said he kicked and punched Mr.~Li, sending him headlong into a doorframe. Mr.~Li died four days
later.

But if the authorities thought they could quell public cynicism with newfound openness, they were
disappointed. Soon after disclosing the identities of the ``volunteers,'' Web users investigated
their backgrounds, revealing that nearly all the ``randomly selected'' investigators were current or
former employees of the state-run media. The team leader, Zhai Li, had previously worked as an
``Internet commentator,'' a euphemism for those who seek to shape public debate with pro-government
postings.

Over the past week, more than 70,000 postings accumulated on QQ.com, one of China's most popular
bulletin boards, and many of the comments were less than sympathetic. ``This kind of hide-and-seek
investigation looks like an indignity to justice,'' read one typical posting.

Even The Beijing News criticized the process as doomed from the start. ``We've learned two lessons
from this case,'' it said. ``One is that such inspections cannot solely rely on passion and idealism
but on rational and pragmatic work. The second is that this has again exposed the flaws of the
current system and that the need to establish an independent investigation system involving experts
is urgent.''

In the report they issued on Monday, committee members acknowledged that their inquiry had been
hamstrung by officials claiming that it might compromise the police investigation. ``We suddenly
recognized that while netizens may be all powerful on the Internet, they are very much helpless in
real life,'' said Mr.~Zhai, the team leader.

On Tuesday, Gong Fei, a provincial official, defended what had become known as the ``elude the cat''
episode, saying it represented a new chapter in grass-roots activism. ``In the future we might still
organize this kind of group to get involved in incidents, if allowed by law,'' he said. ``We hope
the truth of all incidents in the future will be promptly and accurately released so no one has to
be left guessing.''

His exuberance, however, is apparently not widely shared. A poll sponsored by Sina.com, one of
China's top Web portals, showed that as of Tuesday, 86 percent of respondents had said they still
did not believe the police version of Mr.~Li's death.

\section{Preparing for a Flood of Energy Efficiency Spending}

KNOXVILLE, Tenn. -- To the casual eye, the basement of this city's Firehouse 9 looks like a jumble
of old hydrants, Dr Pepper cartons, rakes and random gear. To specialists in energy efficiency, the
1960s-era building is a mess of a different sort: wasteful hot water heaters for the firefighters'
showers, ancient refrigerators and outdated lights.

Wrapping up an elaborate energy audit, Knoxville is about to find out which of 99 city buildings are
wasting the most energy. It hopes to begin repairs this summer, just in time to catch a tsunami of
federal stimulus money earmarked for such unglamorous tasks as replacing light bulbs and fixing
leaky insulation.

Knoxville's timing is excellent. The city began the arduous work of cataloging deficiencies before
the stimulus bill passed, and it is well along in planning its next steps. But experts worry that
other beneficiaries, especially cities, are not ready to oversee the huge sums of energy-efficiency
money about to come their way.

The money in the bill is enough to pay for a tremendous expansion of efficiency efforts across the
country. But as with other parts of the stimulus package, the efficiency plan is creating tension
between spending the money quickly, to get rapid economic stimulus, and spending it well, to do the
most good over the long run.

``There's enormous opportunity here for expansion of energy efficiency in this country,'' said
Lowell Ungar, the policy director for the Alliance to Save Energy, an advocacy group. ``But there is
certainly the potential for waste.''

President Obama signed the stimulus package into law on Feb. 17, hailing it as a shot of money big
enough to help shake the economy from its lethargy while advancing many of his campaign priorities.
Accelerating the country's energy transition is at the top of his list. Many experts in the field
agree with him that carefully chosen investments in efficiency will ultimately save more than they
cost, by cutting energy bills.

At least \$20 billion in the stimulus bill was earmarked for programs like improving the efficiency
of government buildings and the homes of poor people, and trying to find better ways to save energy.
That is far more, advocates say, than any bill in history. Within a few months, the money is likely
to start landing in the bank accounts of thinly staffed state and city agencies that are accustomed
to scraping for a dime here, a dollar there.

Utah expects that its state energy office will receive \$40 million for energy efficiency, renewable
energy and related programs -- 123 times the size of the office's current budget, said Jason Berry,
who manages the four-person unit. He is about to go on a hiring spree.

The package contains \$5 billion to weatherize low-income homes through the Department of Energy,
enough to give the state programs that manage that work 10 to 30 times the money they received last
year, said Christina Kielich, a department spokeswoman.

For advocates of this relatively obscure program, ``it's like they finally got to the other side of
the desert and it's pouring rain,'' said Seth Kaplan, a vice president of the Conservation Law
Foundation, an environmental group.

The stimulus package also contains \$4.5 billion to modernize federal buildings and \$2.5 billion
for research into energy efficiency and renewable energy. The biggest chunk, \$6.3 billion, will be
distributed by the Energy Department in grants to state and local governments, which can spend the
money on things as diverse as thicker window panes for state capitols and rebates for homeowners who
change their light bulbs.

Homes and commercial buildings account for 39 percent of national energy consumption. Experts say
that improving their efficiency is not only cost-effective but also a good way to reduce the
nation's emissions of the greenhouse gases that cause global warming.

But figuring out how to spend the money effectively -- learning which university buildings need
their doors caulked, for example, or which firehouse walls have insulation that is too thin -- can
involve time-consuming, tricky analysis by skilled technicians.

``People are very conservative about their buildings,'' said Donald Gilligan, the president of the
National Association of Energy Service Companies, a trade group. ``Nobody wants to put a failed
technology into the school buildings or have the lights not work.''

In Knoxville, a team of auditors hired by the city is spending six months peering into the grimy
nooks of fire and police stations and even the convention center, where one employee referred to the
downstairs boiler area as a ``money-eating room.''

Knoxville -- which says the stimulus money may help accelerate or expand its program -- hopes to
reduce the city's energy bills as much as 25 percent, and the city is ``definitely on the front end
of the wave as far as efficiency and municipalities addressing efficiency,'' said John Plack Jr., a
director of project development for Ameresco, which is conducting the Knoxville energy audit.

In the Southeastern region of the country, where Mr.~Plack works, low electricity prices have often
made saving energy an afterthought, unlike in California and much of the Northeast. For example,
Nashville, nearly 200 miles west of Knoxville, has not conducted an energy audit of its city
buildings, though it hopes to use stimulus money to look through its own stock of fire stations and
libraries.

``There's a lot of municipalities out there who are completely unaware this is moving forward,''
Mr.~Kaplan said, referring especially to smaller cities. ``They just don't have the infrastructure
in place to deal with this.''

The Energy Department, which is doling out most of the grants, has been assailed on Capitol Hill for
delays in disbursing other types of assistance for clean energy. Ms.~Kielich said in an e-mail
message that the department hoped efficiency grants would begin flowing to city and state energy
offices within 120 days, and that it planned to begin disbursing weatherization money
``expeditiously and responsibly.''

On the receiving end, absorbing the huge increase in money for weatherization could be particularly
challenging, said Ian Bowles, the secretary of energy and environmental affairs for Massachusetts.
Though he contends it can be done, ``the weatherization folks are going to have to quintuple their
effort in order to put that money out,'' he said.

In some cases, the managers of efficiency programs may not need to look far to find ways to spend
the money.

In Knoxville, the Community Action Committee, whose operations include helping poor people
weatherize their homes, works from a building with a \$14,000 monthly utility bill -- some of it
because of an enormous skylight that lets in too much blistering Tennessee sunshine in the summer.

``It's embarrassing,'' said Barbara Kelly, executive director of the committee. ``We do better for
our clients than we do for us.''



\section{A Clinton Listening Tour, but China Gets an Earful}

BEIJING -- Hillary Rodham Clinton said she wanted to spend her first trip as secretary of state
listening. But she ended up saying a lot, and in blunter terms, than many of her pinstriped
predecessors.

The Clinton Straight Talk Express made its last big stop in Beijing on Sunday, with Mrs.~Clinton
explaining to a Chinese talk show host why China had better keep buying United States Treasury
bonds.

``It's a good investment, it's a safe investment,'' she told Yang Lan, the host of a show on Dragon
Television.

The Chinese government, she said, has an even more compelling incentive to keep buying: it needs the
United States to recover as a market for Chinese goods. To jolt the economy back to life, she added,
the United States needs to be able to take on more debt.

``We are truly going to rise or fall together,'' she said. ``We are in the same boat and,
thankfully, we are rowing in the same direction.'' Treasury Secretary Timothy F. Geithner could not
have said it better.

Speaking of Mr.~Geithner, Mrs.~Clinton cleared up the question of how the two of them were going to
divide China policy. Speculation about who would control China policy has crackled since it became
clear she was not ready to cede it to the Treasury Department, as happened during the Bush
administration.

Mrs.~Clinton said she and Mr.~Geithner would share leadership of a high-level consultation between
China and the United States that would be balanced between strategic issues like North Korea and
economic concerns.

Before leaving Beijing, Mrs.~Clinton met with female lawyers, doctors and academics. It is the kind
of encounter she has thrived on since she was first lady, and she has been a heroine in this circle
since her famous speech to a women's conference in Beijing in 1995.

One woman spoke of posing with a wax figure of Mrs.~Clinton on a trip to the United States. Several
had met her before, and they could not resist commiserating with her about her presidential
campaign.

Xie Lihua, the founder of a group that represents rural women, said she memorized the words ``I hope
you win'' for a meeting with Mrs.~Clinton last year. When she shook her hand then, Ms.~Xie was so
nervous that she blanked on the words. ``If I had said it, you probably would have won,'' she said.

In Obama's Town

There was little chance of Mrs.~Clinton's forgetting the campaign when she landed in Jakarta,
Indonesia's capital. Waiting for her on the tarmac was a group of children from the school President
Obama attended when he lived in Indonesia as a boy. They serenaded her under a stormy tropical sky.

After a meeting with the foreign minister, Hassan Wirajuda, the subject of when Mr.~Obama might
visit came up. ``We cannot wait too long, and I wish Hillary Clinton would convey this to President
Obama,'' he said.

At dinner on Wednesday night, Mrs.~Clinton joked about how many people had asked her the same
question. The next morning, she said she had spoken to Mr.~Obama by phone and passed along the
requests. The word is Mr.~Obama might make it to Indonesia in the fall.

Tea With the Empress

After 16 years on the world stage, there is little in the way of pomp and ceremony that can dazzle
Mrs.~Clinton. In Tokyo, she was invited to a formal tea with Empress Michiko, who accorded her the
very rare honor of emerging from the cloistered imperial residence to greet her.

Clasping Mrs.~Clinton's left arm and drawing her close, Empress Michiko spoke in English to her
guest, her eyes sparkling. As they went inside, photographers were ordered to stop their clicking,
one of many rules that keep the imperial family at a strict remove from ordinary society.

Yet at a town hall meeting at the University of Tokyo, Mrs.~Clinton talked about the visit as though
it had been just a couple of old pals catching up. Noting that she first met the empress 15 years
ago, Mrs.~Clinton said, ``We were both saying we had gotten older.'' This, she said, led to a
discussion of how caretaking for the elderly could become a growth industry in Japan.

A Familiar Feel in the Air

There were times when the secretary's plane felt more like a campaign charter, hopping between a
state fair in Jakarta and a stump speech in Seoul, South Korea. Mrs.~Clinton brought along a coterie
of aides from her campaign and the White House, giving the diplomatic mission a distinctly political
flavor.

Lissa Muscatine, a longtime aide, was busy writing speeches. Kiki McLean, a veteran Democratic Party
strategist who is helping Mrs.~Clinton on a transitional basis, trooped back to the press section to
chat up reporters and dole out good-natured spin.

Huma Abedin, Mrs.~Clinton's aide-de-camp and ubiquitous seatmate during the campaign, was never more
than a few steps away from her boss. Ms.~Abedin even accompanied Mrs.~Clinton to her tea with
Empress Michiko, though she was guided to an antechamber to wait.

Even the itinerary had the flavor of a campaign. On the way home to Washington, Mrs.~Clinton's plane
landed at Yokota Air Base in Japan to refuel, and she squeezed in a pep talk to about 350 Marines
and airmen.

Speaking in a hangar under a giant American flag, flanked by two helicopters, Mrs.~Clinton reported
on her trip and praised the troops for being ``on the front lines of defense.'' Then she worked the crowd.

\section{Facing Counterfeiting Crackdown, Beijing Vendors Fight Back}

BEIJING -- Any tourist who has stepped foot in this city's famous Silk Street Market can testify
that it is home to some of the wiliest, most tenacious vendors who ever tried to palm off a fake
handbag on a na\"ive foreigner.

So when the market managers temporarily shut down 29 stalls over the past month for selling
counterfeit goods, no one expected the merchants to acquiesce quietly to the loss of business.

``We expected trouble,'' said Zhao Tianying, a legal consultant with IntellecPro, a Beijing firm
specializing in intellectual property rights, who represents five foreign luxury-brand manufacturers
that have sued the market for trademark violations. ``But we never imagined this.''

The vendors have responded with the same ferocity with which they nail down a sale. Dozens of them
have staged weekly protests against IntellecPro lawyers who are pursuing the trademark case, mocking
them as bourgeois puppets of foreigners. The vendors confronted witnesses who provided evidence of
trademark violations and filed a countersuit asserting that only the government can shutter a
business.

A few characters scrawled in pencil on the wall outside IntellecPro's office sums up the vendors'
message: ``We want to eat!''

The skirmish between the crafty but mostly uneducated hawkers and five of the world's best known
producers of designer goods is part of a much bigger fight over China's vast counterfeit industry.
American movie, music and software companies alone estimate that Chinese pirated goods cost them
more than \$2 billion a year in sales.

Any successful product is likely to be illegally copied in China, warns the Web site of the American
Embassy here. China's government has pledged to crack down, and it faces increasing pressure to show
progress. But some doubt much will change until China graduates from manufacturing goods to
designing them, and has more to lose than gain.

The Silk Street Market case suggests that change is slow and painful.

It has been four years since Burberry, Gucci, Chanel, Louis Vuitton and Prada first sued the
market's operator, the Beijing Silk Street Company, and individual vendors for trademark violations.
Only now has the legal pressure produced tangible results.

As part of a court-mediated agreement, the market's managers agreed to punish offending vendors,
shutting down six to eight at a time for up to a week. George Wang, the market's general manager,
said the manufacturers threatened to renew their suit if sales of counterfeits were not curtailed in
six months.

In response, dozens of vendors descended on IntellecPro's office on Feb. 4, occupying the reception
area for hours while the police tried to mediate, said Ms.~Zhao, the legal consultant. The next day,
she said, they stormed past the receptionist, banged on the walls and swore at the staff. The firm's
senior partner, Hu Qi, was afraid to go home and slept in a hotel for three nights.

Last Monday, more than 50 vendors showed up for the sixth protest. They waved signs and chanted
slogans outside the firm's building while IntellecPro lawyers, with 12 hired guards on hand, had
their lunch delivered.

``We are trying to run businesses here,'' said one 37-year-old vendor in a red coat, a fake Dolce \&
Gabbana handbag on her arm. ``They don't have any proof.'' She refused to give her name, saying she
already faced enough scrutiny.

Asked about her handbag, she insisted: ``We don't read English. We don't know what the letters mean.
We just think it is pretty.''

Another vendor, 24, who gave her last name as He, said: ``We want to be compensated for our losses.
And we want a public apology.''

Mr.~Wang, the market's amiable, 43-year-old manager, said he was ``stuck in a terrible position.''

``The five brands are saying, 'You are not doing a good enough job in protecting our intellectual
property rights,' '' he said. ``And the vendors are saying, 'You are going overboard in protecting
intellectual property rights.' But hey, what can we do? We would rather be known in the world as
going overboard than for not.''

There is little risk of that now. Tourist guidebooks call the Silk Street Market, a seven-story
glass box near Beijing's diplomatic quarter, one of China's most popular spots to buy cheap,
good-quality imitations. With some 1,200 stalls, it attracts 15 million shoppers a year, two-thirds
of them foreigners, Mr.~Wang said.

In the noisy basement, hawkers of leather goods buttonhole passing foreigners, cajoling until all
hope of a sale is lost. They chat easily in broken English and can assess a copy's quality in
seconds; the best, rated ``super-A,'' are almost indistinguishable from genuine products.

Their shelves bulge with fake handbags bearing the designs and tags of Coach, Dolce \& Gabbana,
Chlo\'e and other famous companies, which, Mr.~Wang said, ``have not come to us with a complaint.''

Fake Gucci and Louis Vuitton bags are still offered, but are hidden inside cupboards; buyers are
invited to seal the transactions outside the building. ``There is too much pressure on right now,''
whispered one vendor, a few yards from a stall shrouded in a gray curtain.

Xu Shengzhong, the vendors' lawyer, tries to portray his clients as too ignorant to distinguish fake
goods from real or to recognize brand names. ``They have no idea this says Louis Vuitton,'' he said,
tapping a brown wallet with the brand's distinctive logo.

More potent than such equivocations, though, may be the threat of social unrest. The police, saying
their first priority is to maintain order, organized a meeting between some vendors and Mr.~Hu,
IntellecPro's senior partner, according to the firm's spokeswoman. It went poorly.

The merchants lectured Mr.~Hu on the need for intellectuals like him to respect workers, while
Mr.~Hu tried to defend his patriotism, according to an audio CD made from one vendor's clandestine
cellphone recording.

``These ordinary people work for decades, to their deaths!'' one vendor said. ``How can you say you
are patriotic?''

Mr.~Wang said that while he sympathized with the vendors, ``the Silk Market must fundamentally
change'' and shift its focus from counterfeit goods to genuine pearls, silks, homegrown brands and
tailoring services.

Last year, the market began its own line of products, warning counterfeiters to stay clear.

Mr.~Wang said he hoped that shoppers changed their habits, too. At present, ``they want the
knockoffs,'' he said ruefully. ``You can see it in their eyes. That is the brutal reality.''

\section{China Outlines Ambitious Plan for Stimulus}

BEIJING -- Warning that China faces ``unprecedented difficulties and challenges,'' Prime Minister
Wen Jiabao outlined a barrage of construction, increased subsidies and economic measures on Thursday
aimed at continuing his nation's modernization despite a world financial crisis.

He also indicated that China's leaders would seek to begin a fundamental shift in their economic
strategy by encouraging citizens to spend and consume more goods, as in most Western economies.
China's startling growth has been driven so far by exports and abundant spending on roads, dams and
other infrastructure projects, a trend that experts say cannot be sustained in the long term.

But in a long speech to the National People's Congress, China's legislature, he did not explicitly
announce any new spending to combat the financial crisis beyond the \$585 billion China committed to
spend in November.

Mr.~Wen's speech, which cited numerous budget increases for construction and other programs, left it
unclear whether those figures were part of the already-announced stimulus program or represented
entirely new spending. Details of the stimulus package have remained vague since its unveiling, and
Thursday's speech shed little light on them.

China's stock market rose about 6 percent on Wednesday, and other global markets also advanced, on
speculation that Mr.~Wen would announce more economic stimulus spending.

Instead, he said that while 2009 would be a year of ``arduous tasks,'' China's economic strategy was
sound and the nation's fundamentally positive outlook was unchanged.

His address to the legislature's 3,000 delegates, which was similar to an American State of the
Union speech, laid out an economic program that analysts here say was dominated by two concerns:
riding out the global downturn and keeping citizens' unhappiness with their lot from boiling over
into public unrest.

``Social stability is definitely of concern to the government; it's always been the top concern,''
Tao Wang, a Beijing-based economist for UBS Securities, said in an e-mail message. ``Everything from
delivering growth to holding down unemployment, it's all about stability.''

Ms.~Wang said that the stimulus program announced in November, which focuses mostly on major
construction projects, was a good way to combat China's flagging economy quickly. But social
investments of the kind Mr.~Wen announced on Thursday, she said, are equally important for China's
economic health.

In his address, Mr.~Wen set a target of 8 percent growth in China's gross domestic product in 2009,
the same as in previous years. The government has long said that this rate is needed to hold down
unemployment and the potential for unrest. The economy expanded 9 percent last year, even after
growth slowed sharply in the last quarter.

But a number of economic experts believe that 6.5 or 7 percent growth, meager by recent Chinese
standards, is increasingly likely.

Mr.~Wen said that the central government would significantly increase spending on schools, hospitals
and clinics, low-income housing, environmental programs and other projects aimed at improving
people's lives.

He also reaffirmed plans to raise subsidies to farmers, old-age pensions and income grants to
China's poorest citizens, and said spending on ``social safety net'' programs would jump 17.6
percent, or about \$6.4 billion.

But he reserved some of his strongest words for the Chinese consumer, whose legendary frugality has
kept China's consumer goods and services industries from growing at anywhere near the rate of the
national economy. Mr.~Wen said bolstering consumer demand must become ``a long-term strategic
principle and a basic point of departure for stimulating economic growth.''

Persuading Chinese consumers to abandon longstanding savings habits may be difficult. Official
documents released on Thursday indicated that the government planned to increase consumer loans,
reduce some V.A.T. and income taxes, cultivate markets for new and used homes and cars, and
subsidize the purchase by rural residents of consumer appliances, like water heaters and air
conditioners.

Spending on an array of major programs, including highway and railroad building, low-income housing
and energy conservation, would more than double, to about \$132 billion.

Such spending and tax cuts, and a general reduction in tax revenues because of a slowing economy,
will push this year's budget deficit to about 3 percent of G.D.P., a modern record.

China's parliamentary session is generally dominated by the pro forma passage of laws and
regulations decided well in advance. The delegates, the handpicked cream of the Communist Party
establishment, have little taste for revolt. They gathered in the Great Hall of the People, on
Tiananmen Square, for Mr.~Wen's overview of domestic and foreign policy goals. Much of the rest of
the session is to be spent giving those goals ritual approval.

In advance of Mr.~Wen's address, officials announced Wednesday that they would raise military
spending by 15 percent this year, continuing a long streak of double-digit jumps that has alarmed
the United States. The rate is slightly lower than in the past few years, however.

The government is eager to use the session to present a united and confident front to a public
battered by rising unemployment and falling incomes. Officials estimated a month ago that 20 million
of the nation's 130 million migrant workers -- the cheap-labor force that powered much of China's
construction and export booms -- had lost their jobs.

A signal goal of China's stimulus program is to ensure that the idle jobless do not become an engine
of social unrest.

Experts will be watching Mr.~Wen's proposals closely to see not only how the proposed \$585 billion
(or 4 trillion yuan) stimulus will be spent, but how much is genuinely new spending. China's central
government is allotting only about \$175 billion of the total; the rest is supposed to come from
banks, investors and local governments, whose finances are especially opaque.

Most of the money is set to go to infrastructure projects, like roads and dams, that pump money
quickly into the economy. But many outside experts and some party figures are calling for more money
to be spent directly on people's basic needs.

\section{Drug Investors Lose Patience}

As merger mania plays out among the pharmaceutical giants, a different sort of financial frenzy has
seized some small, struggling drug makers. Investors are demanding that stragglers close up shop and
hand over any remaining cash.

That is what happened to one company, Avigen, after its most promising drug failed in a clinical
trial last October. Avigen said it would do what countless other biotechnology companies had done in
similar circumstances: move on to the next product in its pipeline.

Not so fast, said its biggest shareholder, the Biotechnology Value Fund. The fund demanded that
Avigen, after 16 years of trial and error, immediately liquidate itself and return its remaining
cash to shareholders.

So much for the traditional model of patience in biotechnology investing, in which companies may
burn through more than a decade and hundreds of millions of venture capital or shareholder dollars
before reaching profitability -- if they ever get there. Now, with cash scarce, credit tight and big
drug companies like Merck intent on branching into biotechnology themselves, struggling start-ups
may no longer get second and third chances to succeed.

In at least eight cases in the last year, anxious investors have tried to block an unsuccessful
biotech company's quest for the next blockbuster, and have fought with management for control of the
corporate carcass. The investors argue that the remaining cash belongs to them and that they -- not
a losing company's executives -- should decide how to invest it.

Some companies, including Avigen, are fighting back. ``I hear that argument'' about shareholder
rights, said Kenneth G. Chahine, Avigen's chief. ``But it's really 'I want to raid the cash.' We're
back to 1987 and 'Barbarians at the Gate.' ''

Such battles have become much more common in recent months, as the stock market crash has pounded
the value of many biotech companies to less than the cash on hand. When that happens, investors can
realize an immediate return if the company dissolves itself -- even if some of the cash will be
consumed in closing the company.

In some cases, investors are succeeding. Under pressure from the hedge fund RA Capital Management,
for example, Northstar Neuroscience, a medical device company in Seattle whose stroke treatment
failed, is proposing to liquidate, with shareholders receiving an estimated \$1.90 to \$2.10 a share
in cash. The company's stock, which had been as low as 90 cents in November, closed at \$1.90 on
Monday.

Another company, Trimeris, whose only product, the AIDS drug Fuzeon, has lost sales to newer
competitors, halted research and development last year and repaid \$55 million -- or \$2.50 a share
-- to stockholders. The company continues in business, but with few employees.

And two companies, VaxGen and NitroMed, have canceled planned reverse mergers because of shareholder
opposition. In a reverse merger, a publicly traded company essentially cedes its cash and stock
listing to a private company with presumably better prospects.

For every Gilead Sciences, which spent \$450 million over 15 years and abandoned its original
technology before becoming profitable, there have been countless ``zombies'' -- companies that lurch
from product to product, surviving years or even decades without ever achieving success.

One company so tarred, by one of its biggest investors, is Penwest Pharmaceuticals.

``The company's history is an unfortunate progression of failed development programs,'' Perceptive
Advisors, an investor in the company, wrote in November to Penwest's board. Perceptive demanded that
Penwest cease all research and development and become a virtual company that would just collect
royalties on its one successful drug. Penwest defended its track record and said it was sticking to
its course.

Some investors say that with capital markets now so tight, the walking dead should be buried to free
up financing for more viable companies. ``It's in a time like this that the good companies are being
dragged down by the bad ones,'' said Oleg Nodelman, a portfolio manager at the Biotechnology Value
Fund.

In some cases, however, the investors asking for their money back are not long-suffering
shareholders. They are speculators who bought in only after the stock price collapsed, hoping to
make a quick killing.

Tang Capital Partners, for instance, began accumulating its 14.9 percent stake in Vanda
Pharmaceuticals only after the Food and Drug Administration rejected Vanda's schizophrenia drug in
July. Tang is now pressing for the company to cease all operations and return cash to shareholders.
Vanda's stock is trading at 80 cents, well below the \$1.74 a share in cash it had as of Dec. 31.

Vanda says that it is still hopeful that it can get its drug approved and that liquidation is not in
the interest of all shareholders.

The Biotechnology Value Fund, often called BVF, was a longtime shareholder in Avigen. But it sold
640,000 shares, nearly all its holdings, for about \$3.95 to \$4.60 a share. The sale was near the
stock's highs for the year -- in the two months before Avigen was scheduled to announce, in October,
the clinical trial results of its drug to treat a symptom of multiple sclerosis.

After the drug failed, BVF swooped in and bought more than eight million shares, nearly a 30 percent
stake, at about 58 cents a share. That was well below Avigen's cash total of about \$1.90 a share at
the time.

BVF has made a \$1-a-share tender offer for Avigen and is trying to replace the directors. If it
gains control, it could liquidate Avigen or sell it to MediciNova, which has said it wants to buy
it. Mr.~Chahine, the chief of Avigen, which is based in Alameda, Calif., said its assets might be
parlayed into a deal that would be worth more than BVF or MediciNova would pay and more than the
liquidation value. ``All we're saying is, give us an opportunity to canvass the field, see what's
out there and bring something to the shareholders,'' he said.

But Mr.~Nodelman said such a process might eat up the company's remaining cash. ``Someone's got to
police the space,'' he said. ``We're making sure that the last \$50 million in the company don't go
to the bankers and the consultants and the golden parachutes.''

BVF, which specializes in smaller biotech companies, has become the most outspoken investor pressing
for its money back. The fund, based in San Francisco, gets about half of its capital from the Ziff
family, which made its fortune in magazine publishing.

Mr.~Nodelman makes no apologies for BVF's having bought Avigen stock again after the collapse. The
fund is also pressing for a cash-out to shareholders from CombinatoRx. BVF has been a continuous
shareholder in the company, although it added to its stake after some CombinatoRx clinical trials
failed.

CombinatoRx, whose strategy is to combine two old drugs to make one new one, has lost \$236 million
since its inception in 2000. The company has about \$1.45 a share in cash, but its stock is trading
for only 66 cents.

Alexis Borisy, the chief executive, said the company, based in Cambridge, Mass., was not ready for
the grave. ``We obviously think there's a lot of upside value in the CombinatoRx technology,'' he
said.

BVF and CombinatoRx are now in confidential discussions about the company's future.

BVF is also one of four investors, which collectively own about two-thirds of the shares, demanding
money back from Neurobiological Technologies of Emeryville, Calif.

The company's stroke drug is derived from the venom of the Malayan pit viper. Three of the
investors, including BVF, were shareholders when that drug failed in a clinical trial in December.
The fourth bought in after the failure. The stock now trades at 58 cents, but its liquidation value
would be as high as \$1 a share.

Matthew Loar, the chief financial officer, said the company was sympathetic to the requests but had
not yet decided what to do. In any case, he said, it could not act as fast as the investors want.

``You can't just turn off the lights in a company in a day,'' he said. Among other things, the
company must figure out what to do with 1,000 poisonous snakes, he said. ``We're going to get rid of
them in the most expeditious, reasonable way possible.''

\section{China's Leaders See a Calendar Full of Trouble}

BEIJING -- Officially, Tuesday is an ordinary day here. Ordinary or not, however, the police will be
on high alert, and some villages across Tibetan areas of western China will be under virtual
military lockdown. For it was 50 years ago on Tuesday that 300,000 Tibetans surrounded the Dalai
Lama's palace in Lhasa, the start of a failed uprising against Chinese rule that resonates to this
day.

The Chinese place enormous emphasis on anniversaries and dates, one more reason that 2009 is shaping
up as a very stressful year for the nation's rulers. Besides navigating an economic crisis that has
rendered millions jobless, the government must steer the country through a sheaf of symbolic dates,
some auspicious, some portentous, all -- the leaders seem to think, anyway -- potential triggers for
public unrest.

Many political analysts here may call their concerns overblown. But history suggests that, perhaps,
the leaders could have a point.

This year commemorates not only the 50th anniversary of the Tibetan revolt. June brings the 20th
anniversary of the crackdown on the 1989 pro-democracy demonstrations in Tiananmen Square, protests
that remain the most visible challenge to Communist rule. April is the 10th anniversary of major
protests by the Falun Gong religious sect, which led to thousands of arrests and, in July of that
year, a government ban on the group.

May heralds the 90th anniversary of the May 4 movement, a 1919 student-led protest against imperial
rule that is both a touchstone of Chinese nationalism and historic proof that people can challenge
their rulers.

Finally, Oct. 1 is the 60th anniversary of the creation of the People's Republic of China. The
government plans a major celebration and will be on the lookout for anyone who seeks to spoil it.

The government is often alert for spoilers, of course, as it was for months during the security
clampdown surrounding the Beijing Olympics last summer. But this year's concerns seem especially
stark. In January, an authoritative state magazine, Outlook, issued an unusually blunt warning of
the likelihood of unrest, saying that ``we are entering a peak period for mass incidents.''

On Monday, state media confirmed that security forces had increased patrols in central Tibet. It was
the first official acknowledgment that police and paramilitary forces had flooded areas where ethnic
Tibetans live, creating an unofficial state of martial law. Foreigners have been barred from many
areas. The ruling Communist Party has established a top-level committee to ensure social stability.
One high-level official in the state media says the committee now has branches at every level of
government and is rumored to be led by Xi Jinping, China's vice president and the heir-apparent to
President Hu Jintao.

News reports say the committee's nickname is 6521 -- a reference to the threats to public order
supposedly posed by this year's 60th, 50th, 20th and 10th anniversaries of major events.

To outsiders, this fixation on dates may seem odd, but it has both cultural and political
explanations. Numerology and the calendar have ancient and honored roles in China, Elizabeth Perry,
a China scholar and professor of government at Harvard, said in an e-mail interview. Even today,
while Americans identify most major events by their locations (the Boston tea party, the Alamo,
Gettysburg, Oklahoma City), the Chinese tend to rely on dates.

The May 4 movement, which mixed anti-imperialism with a call for democracy and human rights, is best
known to most Chinese as 5/4. The 1989 Tiananmen Square protests and crackdown are universally
rendered as the shorthand 6/4.

Ideology also weighs heavily on the penchant for anniversaries, for the Chinese brand of government
propaganda is copied from the Soviet Union's handbook, which whipped even insignificant holidays
into excuses to glorify state power.

While the government more or less ignores the 50th anniversary of the Tibetan uprising, its
rubber-stamp Tibetan Parliament this year declared March 28 to be Serfs Emancipation Day, a
celebration of Tibet's liberation from the Dalai Lama's control.

In that and other respects, Professor Perry wrote, the government's elevation of anniversaries into
politically freighted events runs the risk of backfiring. The bigger the celebration, or the more
galling to dissidents, the greater the likelihood of a reaction.

She said: ``I would not discount the degree to which such anniversaries (like portents of a change
in the mandate of heavens, such as earthquakes, floods and large-scale riots or rebellion) are also
regarded quite seriously by the citizenry at large. There is ample precedent in the republican as
well as communist periods for Chinese protesters to turn the commemoration of political
anniversaries into demands for political change.''

The most recent precedent is Tibet itself, where tensions last year over the 49th anniversary of the
1959 uprising boiled over into violence. But there are many others; the May 4 movement, which was
led by student boycotts in Shanghai and Beijing, has been the inspiration for a number of other
student protests in modern China, including the pro-democracy protests of 1989.

This year's security precautions are so extensive that some analysts discount the chance of any
widespread unrest. And the lot of some Chinese has improved so markedly in the last 20 years,
especially in cities, that the forces that feed protest may be more subdued now, Yuan Weishi, a
professor of philosophy at Sun Yat-sen University in Guangzhou, said in a telephone interview.

After 30 years of Chinese modernization, with its ``traumatic, drastic change every day,'' he said,
average Chinese ``don't want to destroy the social stability they are enjoying now.''

But like some other analysts, he would never say ``never.'' ``People like to take the opportunity of
anniversaries to have their voices heard by the government, to promote reforms in the direction that
people want,'' he said. ``It's possible a small number of people might want to take radical actions
this year. I think that's why the government is paying attention to this.''

\section{Computer Makers Prepare to Stake Bigger Claim in Phones}

The computer industry has hit upon its Next Big Thing. It is called a phone.

Emboldened by Apple's success with its iPhone, many PC makers and chip companies are charging into
the mobile-phone business, promising new devices that can pack the horsepower of standard computers
into palm-size packages.

The companies are also shifting gears because their technological feats of the last two decades --
smaller laptops with faster chips to deliver snazzier graphics -- no longer impress consumers, who
increasingly find their three-year-old computers adequate for everyday tasks.

``The action is really with the smartphones where everyone is competing to cram the most features
into a phone,'' said Linley Gwennap, a veteran chip industry analyst and head of the Linley Group.
``I think of PCs as just kind of boring these days.''

The new smartphones promised by PC companies will, among other things, handle the full glory of the
Internet, power two-way video conferences, and stream high-definition movies to your TV.

It is a development that spells serious competition for established cellphone makers and phone
companies. Apple was the first to spot a sleepy industry, shaking up the handset category two years
ago with the iPhone. Until recently, the handset makers were the ones reacting to the iPhone -- and
then with me-too products.

Now fellow PC makers are announcing plans for smartphones in a variety of sizes, shapes and
abilities.

Acer, the big PC manufacturer, has gone from offering no cellphones to selling eight new models,
with more to come this year.

``The smartphone market is the natural direction of our long-term mobile strategy,'' Gianfranco
Lanci, chief executive of Acer, a Taiwan-based company, said as he announced the products at last
month's World Mobile Conference in Barcelona. ``We're just taking on another dimension.''

Dell has also worked on prototype phones but has not committed to making a new product. And Asustek,
the company that was first to market ultraportable laptops known as netbooks, has new smartphones
coming.

The suppliers to the PC industry have also started shifting to the new market. Intel announced a
deal to supply the cellphone maker LG with chips for new mobile devices. Nvidia, the PC
graphics-chip titan, signed a deal to provide three smartphone makers -- which supply handsets to
brand-name manufacturers and carriers -- with its new Tegra processor.

``The rise of the smartphone and things like graphics and 3D images weren't important when the
incumbents built this business,'' said Michael Rayfield, the general manager of Nvidia's mobile
business unit. ``This is a once in a lifetime deal where a huge market changes the things that are
important to it.''

With smartphones and PCs taking on many of the same functions, there is certainly a fear among PC
makers that if they do not get into cellphones, cellphone makers will start building PCs. Acer has
characterized the smartphone business as a volatile battlefield, saying it needs to fire first and
go after the cellphone makers before they come after it. Indeed, Nokia, the world's largest
cellphone maker, has said it is weighing whether to get into the PC business.

The convergence of the two devices has long been predicted, but it took a confluence of industry
changes for it to begin in earnest. For decades chip manufacturers rushed to leapfrog one another
with faster processors, and computer makers scrambled to squeeze more functions into smaller boxes.
But ever-faster chips eventually become impractical. Their blazing speed requires vast amounts of
power and cooling.

The smartphones give the PC makers a chance to extend their newfound expertise in creating low-power
products.

In particular, Acer hopes to ride its success selling laptops and netbooks into the mobile phone
market through a mix of new software and wireless data plans.

It is working on software that will link all of its portable products together, synchronizing
e-mail, contacts, media files and other information among the products. This could open up a way for
carriers to sell more wireless 3G data services to consumers, since they could offer a single plan
covering multiple devices.

It is an extension of the model that Dell and others are already trying, in which carriers
essentially give \$400 netbooks away to consumers in exchange for two-year contracts to data plans.
Such plans can cost as much as \$1,500 over their lifespan.

Deals that cover laptops as well as phones could prove troublesome for existing cellphone companies,
as it would offer consumers a suite of products that were tightly integrated and supported. In
addition, PC manufacturers come from an industry very familiar with low profit margins and tight
cost structures, and would bring those pressures to bear on established cell manufacturers.

``Acer has learned to live and prosper on very thin margins,'' said Aymar De Lencquesaing, the head
of the company's smart hand-held business group. ``I think we bring this kind of experience to a
market that has perhaps has not yet had to endure similar rigor.''

Both Acer and Nvidia have promised very low-cost smartphones, threatening the most lucrative part of
the cellphone makers' business.

At the same time, the phone market has been bombarded with operating systems from Microsoft, Google
and Intel.

There is a concern among longstanding players in the industry that operating systems and phone
designs are becoming commodities, and that the barrier to entering the marketplace is lower than
when mobile-phone manufacturers were building each handset from scratch.

This gives companies like Motorola and Nokia an entirely new set of problems besides falling sales
and shrinking margins.

``It's cataclysmic for the phone guys, who were used to playing golf on Wednesday afternoons,'' said
Roger Kay, president of Endpoint Technologies Associates, a research firm. ``Those times start to
look pretty good now.''

Not that such a move will be easy for the PC makers. The PC industry has a spotty record for
expanding into consumer electronics. Dell stumbled with its MP3 player, and Hewlett-Packard's line
of televisions failed to catch on with consumers. Both products have been discontinued.

Also, the established mobile-phone makers have longstanding relationships with carriers, which
remain reluctant to provide customer support for a wide array of devices from myriad manufacturers.
Beyond that, traditional cellphone companies do not want to compete with the likes of Microsoft and
Intel, which have grown over the years to dominate the PC business.

Perhaps most critically, traditional phone and mobile chip companies have expertise in making phones
that work.

\section{Yahoo Reverses Its Web Strategy With Web Videos}

For Yahoo's original video strategy, this will be Take Two.

On Monday, the Web portal will announce the latest in a series of niche Web shows. The short
segments about celebrity mothers, titled ``Spotlight to Nightlight,'' are a stark departure from the
company's initial forays into TV-style production for the Internet. This time, Yahoo's executives
say they have found a sustainable model for making original video online, in part by explicitly not
competing with television.

During the middle of the decade, the technology company dreamed up plans for elaborate talk shows,
sitcoms and other TV-type shows. But the expensive attempts to transport TV entertainment to the
Internet ``were all disasters,'' said Trip Chowdhry, a senior analyst for Global Equities Research.
The attempted Hollywood makeover -- which at one point even included plans for an interactive
hidden-camera reality show -- was scrapped in 2006.

Since then, Yahoo has acted as a distributor for audio, video and photographs from other media
companies. And it has continued to produce its own Web shows, albeit quietly, without any of the
fanfare that accompanied its earlier ambitions.

``We may have come at it from the wrong way in the previous era,'' said Sibyl Goldman, the head of
entertainment for Yahoo. Ms.~Goldman said that some of those shows were developed even when they
``didn't target any existing need.''

``The one-way model -- 'We think this is great, we hope you come watch it' -- may have been more of
a TV mindset,'' she said.

Last year, Yahoo ended one of the last holdovers from the Hollywood era, a daily distillation of Web
videos called ``The 9.'' Despite a sponsorship by Pepsi, the show was viewed as an experiment, not a
``response to any specific audience demand,'' said James Pitaro, a vice president who oversees media
for Yahoo.

Now the Yahoo producers are doing the reverse, finding their biggest audiences and then building
short Web shows for those groups of people. Instead of producing TV, Yahoo now recaps TV in a daily
show called ``Primetime in No Time.'' Yahoo says the two- to five-minute-long show has an average of
400,000 daily streams, making it one of the most popular recurring series made for the Web.

The shift in strategy comes as the Web video market matures and media companies seek profitable
projects in a battered advertising market. Yahoo, which can drive users to the videos from its
popular home page, has signed up long-term advertising sponsors for each of its original shows; the
TV recaps are sponsored by Verizon Wireless and the celebrity mother segments will be supported by
State Farm Insurance.

State Farm had been looking for new ways to reach a female audience when it contacted Yahoo about a
partnership, said Ed Gold, an advertising director for the insurance company. Some viewers may treat
the ``Spotlight to Nightlight'' videos as disposable -- nannies and zany baby names are among the
topics -- but Mr.~Gold said they allow State Farm to ``own the environment around the discussion.''

The segments will be hosted by Ali Landry, a former Miss USA. In a twist, the videos are being
produced in both English and Spanish, owing in part to Ms.~Landry's fluency in both languages.

``Spotlight to Nightlight'' was a natural extension of OMG, Yahoo's upbeat entertainment project.
Ms.~Goldman said the producers of the nearly two-year-old site saw traffic spikes around stories
about celebrity baby names and photos and ``started programming more content'' that related to the
subject.

The ability to mine search queries and traffic data to better identify user interests is an
advantage for companies like Yahoo, Mr.~Chowdhry said. Besides TV and celebrity mothers, Yahoo also
produces technology news videos for its finance Web site and game highlights for its sports site.
While they may lack Hollywood ambition, they also lack risk for Yahoo's bottom line, a plus for the
company as it embarks on yet another reorganization.

``These are shows that identified needs,'' Mr.~Pitaro said.

\section{Government Combs Through G.M.'s Survival Plan}

WASHINGTON -- Treasury officials and management experts hired by the Obama administration quietly
began combing through General Motors' latest downsizing plan in Detroit last week, in a last-minute
effort to assess whether more government aid could make the company viable, or whether the better
choice was a managed bankruptcy.

President Obama faces a deadline of March 31 to decide the fate of G.M., and by extension its huge
network of suppliers. In interviews, however, administration officials said they would not be bound
by that date, when Mr.~Obama is scheduled to visit London for a summit meeting on the global
economic crisis.

During the weekend, administration officials said they expected the team of about two dozen Treasury
officials to continue working in Detroit this week, even while G.M.'s two top executives traveled to
Washington to meet the task force that is assembling Mr.~Obama's options.

To judge whether G.M.'s reorganization plan is sufficient, the team has been delving into every
aspect of the company's management, including its somewhat upbeat estimate of how the car market
will look when the recession ends -- a market it projects at over 15 million cars annually, much
higher than Chrysler's forecast — as well as its designs for new products and its financial
controls.

Administration officials said Sunday that the main standard they were using to measure the viability
of G.M. was the probability of recovering additional taxpayer money used to help the company, should
more aid be necessary. Asked about the mission of the team, a White House spokesman, Bill Burton,
said: ``The auto team is working to assess long-term viability within the context of the existing
loan agreements, while recognizing that the economy has deteriorated since September 2008 and that
both companies are now asking for more money.''

While administration officials say that no final decisions have been made, they appear to be giving
more consideration than they were two months ago to a quick, managed bankruptcy. Mr.~Obama has said
he wants avoid that option, because it could further unnerve investors and radiate out through the
company's network of suppliers and dealers.

But the threat of a bankruptcy filing may be the only way to force concessions from bondholders, who
hold \$27 billion of the company's unsecured debt. So far, the bondholders have shown little
willingness to negotiate over a G.M. proposal that would give them a fraction of the face value of
the bonds -- as little as 16 cents on the dollar for some bonds in a deal that could leave them
holding what may turn out to be worthless equity in G.M. The president's task force has added a
bankruptcy lawyer, Matthew A. Feldman of Willkie Farr and Gallagher, to its team, though officials
said that the timing was somewhat coincidental because they had been seeking Mr.~Feldman's help for
some time.

The face-off with the bondholders has quickly become a high-stakes game of chicken. The bondholders
appear to be betting that Mr.~Obama is not willing to take the political risk of letting the company
that once defined the American car industry go bankrupt, a step that could render their bonds
worthless.

The administration has a strong interest in trying to persuade the bondholders' committee that if
they have no other option for reorganizing the company, the government is willing to take that step.
Senior administration officials who met with the United Automobile Workers union last week in
Detroit say that while the meeting was cordial, the union was not ready to give any more concessions
until it sees what happens with the investors who have lent G.M. billions.

``The union is wary of negotiating with themselves,'' said one senior official involved in the
talks. ``They have been asked to give blood four times, and they want to see the bondholders give
some before they do anything more.''

Another senior administration official, who declined to speak on the record because he was not
authorized to talk about the negotiations, described talks with the bondholders as ``the drama that
will play out for the next two weeks,'' as each side waits for the other to blink. ``Is Obama
willing to go the bankruptcy route?'' the official asked. ``He might be, if he thinks it's the best
way to save the most jobs, but we don't know yet.''

That decision may have as much to do with negotiation strategies as it does with economic
fundamentals.

``It seems that the dispute between G.M. and its constituents will not be settled with an economic
solution,'' said John Casesa, whose firm, Casesa Shapiro Group, advises auto companies. ``This is
now in the realm of politics. Is it more politically palatable for G.M. to go into a
government-sponsored bankruptcy than to be sustained by more loans?''

The prospect of the bondholders taking huge losses is critical to the plan that G.M. submitted last
month, in which it aimed to eliminate roughly two-thirds of its unsecured debt by March 31. Its
other commitment was to reduce labor costs by the government-imposed deadline as part of the terms
of the \$13.4 billion in federal loans it already has. But with only two weeks to go, the company
has been unable to get the union and its bondholders to agree on concessions to ensure G.M.'s
long-term viability.

The U.A.W. has agreed to givebacks on job security and work-rule provisions of their national labor
contract. But the union's president, Ron Gettelfinger, has not yielded to pressure to accept G.M.
stock in lieu of cash to finance 50 percent of a new health care trust for 391,000 U.A.W. retirees
and surviving spouses.

G.M. owes more than \$20 billion to the health care trust, which will assume responsibility for
billions of dollars in annual medical bills for U.A.W. retirees beginning next year.
Mr.~Gettelfinger recently reached a deal with the Ford Motor Company to accept company stock to
finance half of Ford's health care trust.

But Mr.~Gettelfinger has vowed not to make a similar deal with G.M. unless the company's bondholders
agree to swap their debt for equity in the company.

The bondholders are being asked to take equity in exchange for reducing G.M.'s \$27 billion in
unsecured debt by two-thirds. So far, a bondholders' committee has resisted the debt-for-equity
exchange, given G.M.'s depressed stock price and its uncertain future.

\section{In Downturn, China Exploits Path to Growth}

GUANGZHOU, China -- The global economic downturn, and efforts to reverse it, will probably make
China an even stronger economic competitor than it was before the crisis.

China, the world's third-largest economy behind the United States and Japan, had already become more
assertive; now it is exploiting its unusual position as a country with piles of cash and a strong
banking system, at a time when many countries have neither, to acquire natural resources and make
new friends.

Last week, China's prime minister, Wen Jiabao, even reminded Washington that as one of the United
States' biggest creditors, China expects Washington to safeguard its investment.

China's leaders are turning economic crisis to competitive advantage, said economic analysts.

The country is using its nearly \$600 billion economic stimulus package to make its companies better
able to compete in markets at home and abroad, to retrain migrant workers on an immense scale and to
rapidly expand subsidies for research and development. Construction has already begun on new
highways and rail lines that are likely to permanently reduce transportation costs.

And while American leaders struggle to revive lending -- in the latest effort with a \$15 billion
program to help small businesses -- Chinese banks lent more in the last three months than in the
preceding 12 months.

``The recent tweaks to the stimulus package indicate a sharper focus on the long-term
competitiveness of Chinese industry,'' said Eswar S. Prasad, a former China division chief at the
International Monetary Fund. ``Higher expenditures on education and research and development, along
with amounts already committed to infrastructure investment, will boost the economy's
productivity.''

The international economic slowdown is also doing some things that Chinese authorities had tried and
failed to do for four years: slow inflation, reverse what had been an ever-growing dependence on
exports and pop a real estate bubble before it could grow even bigger.

The recession in most of the large economies in the world is inflicting real pain here -- causing a
record plunge in Chinese exports, putting 20 million migrant workers from within China out of their
jobs and raising the potential for increased and sustained social unrest. But as President Hu Jintao
told the National People's Congress last week, ``Challenge and opportunity always come together --
under certain conditions, one could be transformed into the other.''

To that end, Chinese companies are shopping for foreign businesses to acquire. The commerce ministry
announced late Monday that it was greatly easing the government approval process for Chinese
companies seeking permission to make foreign acquisitions.

The ministry is now leading its first mergers and acquisitions delegation of corporate executives to
Europe; the executives are looking at companies in the automotive, textiles, food, energy,
machinery, electronics and environmental protection sectors.

The government initiatives coincide with some immediate benefits of the slowdown for China. For
instance, air freight and ocean shipping costs have plunged by as much as two-thirds since last
summer as demand has fallen.

Blue-collar wages, which had doubled in four years in some coastal cities, have fallen for many
workers this winter, causing personal pain but reviving China's advantage in labor costs.

Unemployment has pushed down the piece rates that factories pay for each garment sewn or toy
assembled. Overtime has practically disappeared.

Lao Shu-jen, a migrant worker from Jiangxi province who works at a blue jeans factory here, said
that he earned \$350 a month late last year but would be lucky to earn \$220 a month this spring.

``There are a lot of blue jeans'' piling up in the back of the factory with no sign of buyers, he
said.

Highly qualified middle managers, in acutely short supply a year ago, are now widely available
because of layoffs. They are likely to stay that way -- although white-collar unemployment could
pose a threat of social unrest. Limited job opportunities contributed to the Tiananmen Square
protests 20 years ago.

Some jobs are still available now. Four days after a shoe factory closed here for lack of orders,
laying off several hundred workers, there were four ads on the factory's front gate from other shoe
factories seeking to hire skilled workers.

Unskilled laborers face the greatest difficulty finding jobs. But with subsidies from Beijing,
provincial governments have embarked on large-scale vocational training programs of the sort that
the United States has discussed but not actually tried.

Guangdong province alone, here in southeastern China, is quadrupling its vocational training program
this year to teach four million workers engaged in three-month or six-month programs.

The main comparable program in the United States, under the Workforce Investment Act, has been
training fewer than 250,000 a year, although President Obama's stimulus program provides funding
that could double the number of American workers in training programs.

The Guangdong training programs are half in the classroom and half in the factory, usually the
business that plans to employ the trainees. By increasing productivity, training programs can hold
down corporate labor costs per unit of production for years to come.

China's huge training programs may also help preserve social stability by keeping the unemployed off
the streets, although Chinese officials deny that is their intention.

Multinationals are cutting back less in China than elsewhere -- and some are even expanding.

Intel is shutting down semiconductor production lines sooner than previously planned at older,
smaller operations in Malaysia and the Philippines as it opens a large, new factory in Chengdu in
western China.

IMI Plc., the big British manufacturer of items as diverse as power plant valves and brewery
equipment, has just announced an accelerated shift of operations to China, India and the Czech
Republic, after cutting its global work force by 10 percent since December.

And Hon Hai, the 600,000-employee Taiwanese company that is one of the world's largest contract
manufacturers of products like the Apple iPhone and Nintendo Wii game console, has just increased
employment by nearly 5 percent in China even as it cuts overall employment by 3 to 5 percent.

Yet China's economy still has weaknesses. Little is being done to shift the economy away from a
heavy reliance on capital spending and toward greater consumption. The social safety net of
pensions, health care and education barely exists, so Chinese families save heavily.

Strict government policies on labor and the environment, intended to address serious shortfalls in
both and imposed a year ago when China felt more confident of its economic strength, are prompting
low-tech industries like toy manufacturing to move to other, less stringent countries.

Top labor officials insisted during the National People's Congress that they would resist suggestions from some Chinese executives that the new standards be relaxed.

\section{The Case for Paying Out Bonuses at A.I.G.}

Do we really have to foot the bill for those bonuses at the American International Group?

It sure does sting. A staggering \$165 million -- for employees of a company that nearly took down
the financial system. And heck, we, the taxpayers, own nearly 80 percent of A.I.G.

It doesn't seem fair.

So here is a sobering thought: Maybe we have to swallow hard and pay up, partly for our own good. I
can hear the howls already, so let me explain.

Everyone from President Obama down seems outraged by this. The president suggested on Monday that we
just tear up those bonus contracts. He told the Treasury secretary, Timothy F. Geithner, to use
every legal means to recoup taxpayers' money. Hard to argue there.

``This isn't just a matter of dollars and cents,'' he said. ``It's about our fundamental values.''

On that last issue, lawyers, Wall Street types and compensation consultants agree with the
president. But from their point of view, the ``fundamental value'' in question here is the sanctity
of contracts.

That may strike many people as a bit of convenient legalese, but maybe there is something to it. If
you think this economy is a mess now, imagine what it would look like if the business community
started to worry that the government would start abrogating contracts left and right.

As much as we might want to void those A.I.G. pay contracts, Pearl Meyer, a compensation consultant
at Steven Hall \& Partners, says it would put American business on a worse slippery slope than it
already is. Business agreements of other companies that have taken taxpayer money might fall into
question. Even companies that have not turned to Washington might seize the opportunity to break
inconvenient contracts.

If government officials were to break the contracts, they would be ``breaking a bond,'' Ms.~Meyer
says. ``They are raising a whole new question about the trust and commitment organizations have to
their employees.'' (The auto industry unions are facing a similar issue -- but the big difference is
that there is a negotiation; no one is unilaterally tearing up contracts.)

But what about the commitment to taxpayers? Here is the second, perhaps more sobering thought:
A.I.G. built this bomb, and it may be the only outfit that really knows how to defuse it.

A.I.G. employees concocted complex derivatives that then wormed their way through the global
financial system. If they leave -- the buzz on Wall Street is that some have, and more are ready to
-- they might simply turn around and trade against A.I.G.'s book. Why not? They know how bad it is.
They built it.

So as unpalatable as it seems, taxpayers need to keep some of these brainiacs in their seats, if
only to prevent them from turning against the company. In the end, we may actually be better off if
they can figure out how to unwind these tricky investments.

Not that any of this takes the bite out of paying these bonuses. For better or worse -- in this
case, worse -- someone at A.I.G. decided this company needed to sign bonus agreements last year to
keep people before the full extent of its problems became clear.

Now we can debate why A.I.G. felt it necessary to guarantee seven executives at least \$3 million
apiece when the economy was clearly on shaky ground. Perhaps we will find out these contracts were a
bit of sleight of hand to enrich executives who knew this financial Titanic had hit the iceberg. But
another possible explanation is that A.I.G. knew it needed to keep its people.

That is the explanation offered by Edward M. Liddy, who was installed as A.I.G.'s chief executive
when the government effectively nationalized the company last fall. (He is being paid \$1 a year.)

``We cannot attract and retain the best and brightest talent to lead and staff'' the company ``if
employees believe that their compensation is subject to continued and arbitrary adjustment by the
U.S. Treasury,'' he said.

There's some truth to what Mr.~Liddy is saying. Would you want to work at A.I.G.? Sure, maybe for
\$3 million. But not if you could go somewhere else for even more -- or even much less.

``The jobs are terrible,'' said Robert M. Sedgwick, an executive compensation lawyer at Morrison
Cohen who represents a number of employees of banks that have taken government money. ``You have to
read about yourself in the paper every day. These people are leaving as soon as they can.''

Let them leave, you say. Where would they go, given the troubles in the financial industry? But the
fact is, the real moneymakers in finance always have a place to go. You can bet that someone would
scoop up the talent from A.I.G. and, quite possibly, put it to work -- against taxpayers' interests.

``The word on the street is that A.I.G. employees are being heavily recruited,'' Ms.~Meyer says.

Of course, if taxpayers had not bailed out A.I.G., these contracts would not be worth anything.
Andrew M. Cuomo, the attorney general of New York, made the point on Monday, when he subpoenaed
A.I.G. for the names of the people who received the bonuses. If A.I.G. had spiraled into bankruptcy,
its employees would have had to get in line with other unsecured creditors.

Mr.~Cuomo wants to know who A.I.G.'s lucky employees are, and how they have been doing at their
jobs. So here is a suggestion for him. Get the list, and give those big earners at A.I.G. a
not-so-subtle nudge: Perhaps they will ``volunteer'' to give some of their bonuses back or watch
their names hit the newspapers. But in the meantime, despite how offensive and painful it might be,
let's honor the contracts.

\section{Hadoop, a Free Software Program, Finds Uses Beyond Search}

BURLINGAME, Calif. -- In the span of just a couple of years, Hadoop, a free software program named
after a toy elephant, has taken over some of the world's biggest Web sites. It controls the top
search engines and determines the ads displayed next to the results. It decides what people see on
Yahoo's homepage and finds long-lost friends on Facebook.

It has achieved this by making it easier and cheaper than ever to analyze and access the
unprecedented volumes of data churned out by the Internet. By mapping information spread across
thousands of cheap computers and by creating an easier means for writing analytical queries,
engineers no longer have to solve a grand computer science challenge every time they want to dig
into data. Instead, they simply ask a question.

``It's a breakthrough,'' said Mark Seager, head of advanced computing at the Lawrence Livermore
National Laboratory. ``I think this type of technology will solve a whole new class of problems and
open new services.''

Three top engineers from Google, Yahoo and Facebook, along with a former executive from Oracle, are
betting it will. They announced a start-up Monday called Cloudera, based in Burlingame, Calif., that
will try to bring Hadoop's capabilities to industries as far afield as genomics, retailing and
finance.

The core concepts behind the software were nurtured at Google.

By 2003, Google found it increasingly difficult to ingest and index the entire Internet on a regular
basis. Adding to these woes, Google lacked a relatively easy to use means of analyzing its vast
stores of information to figure out the quality of search results and how people behaved across its
numerous online services.

To address those issues, a pair of Google engineers invented a technology called MapReduce that,
when paired with the intricate file management technology the company uses to index and catalog the
Web, solved the problem.

The MapReduce technology makes it possible to break large sets of data into little chunks, spread
that information across thousands of computers, ask the computers questions and receive cohesive
answers. Google rewrote its entire search index system to take advantage of MapReduce's ability to
analyze all of this information and its ability to keep complex jobs working even when lots of
computers die.

MapReduce represented a couple of breakthroughs. The technology has allowed Google's search software
to run faster on cheaper, less-reliable computers, which means lower capital costs. In addition, it
makes manipulating the data Google collects so much easier that more engineers can hunt for secrets
about how people use the company's technology instead of worrying about keeping computers up and
running.

``It's a really big hammer,'' said Christophe Bisciglia, 28, a former Google engineer and a founder
of Cloudera. ``When you have a really big hammer, everything becomes a nail.''

The technology opened the possibility of asking a question about Google's data -- like what did all
the people search for before they searched for BMW -- and it began ascertaining more and more about
the relationships between groups of Web sites, pictures and documents. In short, Google got smarter.

The MapReduce technology helps do grunt work, too. For example, it grabs huge quantities of images
-- like satellite photos -- from many sources and assembles that information into one picture. The
result is improved versions of products like Google Maps and Google Earth.

Google has kept the inner workings of MapReduce and related file management software a secret, but
it did publish papers on some of the underlying techniques. That bit of information was enough for
Doug Cutting, who had been working as a software consultant, to create his own version of the
technology, called Hadoop. (The name came from his son's plush toy elephant, which has since been
banished to a sock drawer.)

People at Yahoo had read the same papers as Mr.~Cutting, and thought they needed to even the playing
field with their search and advertising competitor. So Yahoo hired Mr.~Cutting and set to work.

``The thinking was if we had a big team of guys, we could really make this rock,'' Mr.~Cutting said.
``Within six months, Hadoop was a critical part of Yahoo and within a year or two it became
supercritical.''

A Hadoop-powered analysis also determines what 300 million people a month see. Yahoo tracks peoples'
behavior to gauge what types of stories and other content they like and tries to alter its homepage
accordingly. Similar software tries to match ads with certain types of stories. And the better the
ad, the more Yahoo can charge for it.

Yahoo is estimated to have spent tens of millions of dollars developing Hadoop, which remains
open-source software that anyone can use or modify.

It then began to spread through Silicon Valley and tech companies beyond.

Microsoft became a Hadoop fan when it bought a start-up called Powerset to improve its search
system. Historically hostile to open-source software, Microsoft nevertheless altered internal
policies to let members of the Powerset team continue developing Hadoop.

``We are realizing that we have real problems to solve that affect businesses, and business
intelligence and data analytics is a huge part of that,'' said Sam Ramji, the senior director of
platform strategy at Microsoft.

Facebook uses it to manage the 40 billion photos it stores. ``It's how Facebook figures out how
closely you are linked to every other person,'' said Jeff Hammerbacher, a former Facebook engineer
and a co-founder of Cloudera.

Eyealike, a start-up, relies on Hadoop for performing facial recognition on photos while Fox
Interactive Media mines data with it. Google and I.B.M. have financed a program to teach Hadoop to
university students.

Autodesk, a maker of design software, used it to create an online catalog of products like sinks,
gutters and toilets to help builders plan projects.. The company looks to make money by tapping
Hadoop for analysis on how popular certain items are and selling that detailed information to
manufacturers.

These types of applications drew the Cloudera founders toward starting a business around Hadoop.

``What if Google decided to sell the ability to do amazing things with data instead of selling
advertising?'' Mr.~Hammerbacher asked.

Mr.~Hammerbacher and Mr.~Bisciglia were joined by Amr Awadallah, a former Yahoo engineer, and
Michael Olson, the company's chief executive, who sold a an open-source software company to Oracle
in 2006.

The company has just released its own version of Hadoop. The software remains free, but Cloudera
hopes to make money selling support and consulting services for the software. It has only a few
customers, but it wants to attract biotech, oil and gas, retail and insurance customers to the idea
of making more out of their information for less.

The executives point out that things like data copies of the human genome, oil reservoirs and sales
data require immense storage systems.

\section{Verizon Customers - Just Say No!}

It is easier to seek forgiveness than it is to get permission according to Verizon, which has once
again shown us what large corporations should not be doing when it comes to customer service.

David Weinberger, co-author of The Cluetrain Manifesto and the more recent Everything is
Miscellaneous received a letter today from Verizon. A ``legalistic pamphlet'' that informed him he
has 45 days to opt out of 'agreeing' to let Verizon share his personal information.

Weinberger, unlike the majority of us who rarely read the associated paraphernalia that arrives with
bills and the like, noticed that Verizon's modus operandi was to share Customer Proprietary Network
Information - the data created as a result of your relationship with Verizon Wireless - unless you
ask them to stop. This information includes ``services purchased (including specific calls you make
and receive), billing info, technical info and location info. They promise to only share this with
'affiliates, agents and parent companies.' It will definitely not be shared with 'unrelated third
parties' \ldots unless, perhaps that third party pays Verizon to become an affiliate, whatever the
heck 'affiliate' means,'' Wienberger wrote. While getting this in the mail is great, what about
those people who only receive electronic copies from Verizon? Did they too receive a message? I
didn't, so I decided to investigate, and sure enough, after logging into my Verizon Wireless account
I couldn't find a mention of the CPNI. As Weinberger pointed out, there is a link to it in my
messages, but alas, as he also mentions, the link is ``not available.''

Verizon has not been having a great run recently. Sure, in January Microsoft opened its wallet to
Verizon and struck a five year partnership with the wireless carrier in hopes of getting access to
its more than 80 million subscribers, but what if those subscribers begin to slowly dwindle away?

Verizon was dealt a blow last month when the US Court of Appeals for the District of Columbia
Circuit denied the company's request to overturn a decision by the Federal Communications Commission
regarding its marketing tactics. Earlier this month the Pennsylvania Public Utility Commission
decided Verizon's customer service ``stinks'' and sought a settlement to improve customer service.
Next week, Delaware customers with complaints about their telecommunications service will have the
opportunity to air them during the states Public Service Commission hearings.

Verizon spokesman Harry Mitchell, in discussing next week's public forum said the company still does
not believe the hearings are necessary. ``There were some issues we did have, but we worked through
those, in collaboration with the commission staff, and made great improvements,'' he said.

While they may have fixed the problems associated with gripes about their old copper pipes in
Delaware, they clearly haven't yet addressed their customer service issues.

For a company that prides itself on operating the nation's most reliable and largest wireless voice
and data network, and spends an absolute fortune on creating masterful ads (YouTube link) to
entertain us, it clearly still has a way to go. Perhaps it's time the company started putting some
of its energy into the Twitter account it set up, and got up close and personal to its community and
their needs.

Note: David Weinberger's post walks you through the process of opting out of Verizon's 'share your
information' default setting. If you haven't done it yet, we highly recommend you do.



\section{Bottling the Magic Behind Google and Facebook}

Cloudera is the quintessential Silicon Valley story.

Three of the top engineers from Google, Yahoo and Facebook have teamed up with an ex-Oracle
executive to tackle the problems inherent in quickly analyzing big piles of data. On Monday, they're
revealing a commercial product based on the open source software Hadoop, which provides the
analytical magic behind the world's biggest Web sites. The team at Cloudera, based in Burlingame,
Calif., think they can extend Web smarts to the business world, aiding companies in retail,
insurance, bio-tech and oil and gas.

Hadoop is the open-source version of the file system and MapReduce technology developed by Google.
Google has used such software to rewire its entire search index, making it possible for the company
to run ever-faster searches on cheap servers and to ask questions of its vast data stores and
receive coherent answers.

Rather than keeping data locked in a central database, Google spreads information across thousands
of servers. Engineers can then send out requests to these servers via MapReduce and gain new
insights into peoples' searching behavior and the relationships between Web sites. Best of all,
MapReduce keeps these complicated jobs humming along even when computers fail because of its ability
to maintain a cohesive picture of all the systems.

While Google has kept the deep details on this technology to itself, the company did publish a
couple of papers describing some of the underlying principles. That gave Doug Cutting, formerly a
software consultant and now a Yahoo engineer, enough information to create an open-source take on
the code.

Yahoo has since invested millions of dollars improving Hadoop and uses the technology to figure out
what users should see on its home page, based on their surfing habits, and what ads to display next
to search results.

Other Web 2.0 users, including Microsoft, Facebook and Fox Interactive Media, have picked up Hadoop
as well.

The founders of Cloudera argue that the analytical powers of Hadoop can benefit a whole new class of
businesses. For example, they want to show biotech firms new ways of analyzing genome and protein
data and give oil and gas firms new ways of digging through their reservoir data.

The pitch has proved attractive enough for Accel Partners to pump money into the start-up. Diane
Greene, the co-founder of VMware; Marten Mickos, the former chief executive of MySQL; and Gideon Yu,
the chief financial officer at Facebook, have invested in the company as well.

While Hadoop remains free, Cloudera plans to sell support and consulting services around the
software.

The backgrounds of the executives point to the classic Silicon Valley nature of the story.

Just 26, Jeff Hammerbacher has already worked on Wall Street and at Facebook after graduating from
Harvard, where he earned a degree in math. Christophe Bisciglia, 28, arrived at Google after raising
and selling horses online during his high school years.

Amr Awadallah, 38, arrived in the United States from Egypt and secured a job at Yahoo, where he
helped develop Hadoop. And Mike Olson, the 46-year-old chief executive, is a database executive that
sold an open-source software maker called Sleepycat to Oracle in 2006.

While Google could make a fuss about intellectual property rights to the technology, the company has
given Cloudera its blessing. Mr.~Bisciglia discussed the company with Google's chief executive, Eric
Schmidt, last March.

``He agreed that this technology is not just for researchers, and it's good for Google to make this
pervasive,'' Mr.~Bisciglia said. ``The more data people create, the more data Google can slurp up.''

A number of prominent computer scientists have hailed Hadoop as the right answer for an age when
companies have moved from dealing with gigabytes of data to terabytes and now petabytes (one
petabyte is equal to 1 million gigabytes or 1,000 terabytes). It's one thing to store all of that
information and another thing to be able to mine it in an efficient manner.

``It is a new reality that people have the ability to store and analyze terabytes and petabytes of
data,'' Bisciglia said. ``Now they need the tools to process it.''

\section{3G Phones Exposing Networks'Last-Gen Technology}

SAN FRANCISCO -- Oh, the things modern mobile phones can do. They are music-playing, video-taking,
direction-providing multimedia powerhouses. But many people have trouble getting them to perform
their most basic functions, like making phone calls.

The underlying problem, industry analysts say, is the complex quilt of the nation's wireless
networks. The major mobile carriers have spent tens of billions of dollars on new voice and data
networks that they advertise as superfast wireless express lanes. But analysts say these upgrades
present major engineering challenges, and the networks often underperform.

The resulting technological glitches have given many owners of fancy new phones the urge to throw
them out the window and onto the highway.

For many, the iPhone has become a symbol of the gap between the promise of a powerful device and the
reality of inconsistent service. Its owners complain of continual hiccups, particularly in certain
cities.

AT\&T, which is the exclusive carrier for the iPhone in the United States, says it has done a lot to
improve its network, and is doing more. Last Tuesday, the company announced plans to invest around
\$11 billion this year to expand and improve its wireless and broadband networks.

``I'm not minimizing the frustration somebody may feel, but I think the improvements in wireless in
this country have been extraordinary,'' said Mark Siegel, a spokesman for AT\&T. ``It's come a long,
long way.''

For some AT\&T customers, more improvements to the company's so-called 3G, or third-generation,
high-speed network can't happen soon enough. And industry analysts say the problems at all carriers
are becoming more glaring as the growing popularity of so-called smartphones puts pressure on their
networks.

``The iPhone and the carriers are shoving 3G down our throats,'' said Edward Snyder, an industry
analyst with Charter Equity Research. ``But the actual experience has been abysmal.''

Overall customer satisfaction with cellphone service has been rising, but it varies among cities and
carriers. Verizon customers tend to be happiest with their service, while AT\&T and Sprint customers
were less satisfied, according to a survey published in Consumer Reports magazine in January. Bob
Goodson, 28, the chief executive of a start-up in San Francisco, upgraded in January from a 2G
iPhone to a 3G version. On the whole he is very pleased, he said, but his experience varies widely
based on what part of the country he is in.

During a recent two-week trip to New York, he said, the coverage was far inferior to what he
experiences in California, and made it tough for him to use the map function to get directions.

``I found myself walking around Manhattan frustrated,'' Mr.~Goodson said. ``It couldn't hang on to
the network.''

Even when the network is within reach, its speeds are often not what they should be. A Gartner
research report released in January found that data speeds for mobile phone users are often half of
what is advertised by the carriers. The most glaring problem, Gartner found, is at AT\&T.

``AT\&T is constantly falling below the threshold,'' said Ken Dulaney, a mobile computing analyst
for Gartner, who said he had heard from three of Gartner's major corporate clients in the last three
months that their employees were frustrated with AT\&T's service. ``I can't say that Verizon is
trouble-free, but we've heard fewer complaints.''

The reasons for the trouble are complicated. Part of the problem is that the companies are
constantly upgrading their networks -- creating a patchwork of technology on cell towers, and
integrating slices of radio spectrum that carry voice and data transmissions.

Analysts said the problem was not unique to AT\&T, but was especially pronounced on its network in
some cities because of the way its infrastructure was built.

AT\&T began introducing its 3G network in 2005, upgrading the equipment and antennas at many of its
40,000 cell towers nationwide. It built the network to complement and take advantage of the
technology servicing its older 2G or second-generation network. Many phones still use the 2G
network, so it must be kept running.

But there are important differences between the 2G and 3G networks, and getting them to work
together presents problems, according to engineers who work on the infrastructure.

Take, for instance, the difference in the way voice and data traffic is carried on the two networks.
On AT\&T's 2G network, cellphone towers -- even ones in close proximity to one another -- use
different chunks of the radio spectrum to carry information. As phone users move around on foot or
in a car, their phones switch from one frequency to another.

On the 3G network, all of the cell towers use the same frequency to transmit information. On its
face, this would seem to make things simpler. But this technology also adds a wrinkle: when phones
get too close to too many 3G towers using the same frequency, they can become overwhelmed with radio
noise.

``When you have more than three cell sites overlapping, you get interference,'' said one
infrastructure engineer who works for AT\&T, who asked not to be named so as not to upset the
company. ``You get bad quality, funky sounds. If you're doing data, the rates get slower and slower
until you lose it.''

Kristin S. Rinne, senior vice president of architecture and planning for AT\&T, said the company had
done a good job of diminishing the prospect of such interference by limiting the strength of signals
from overlapping towers. Moreover, she said, the phones themselves have a role to play; some
handsets, she said, do a good job of managing the interference internally, while others do not.

``The chipset inside the handset impacts how adaptable it is,'' Ms.~Rinne said. She declined to
discuss the iPhone's performance.

The iPhone, of course, is not perfect. Mr.~Dulaney of Gartner said the phone did not communicate
with the data network as efficiently as it could -- which is a problem for a phone whose owners are
among the heaviest users of mobile data.

Greg Joswiak, vice president of iPod and iPhone product marketing at Apple, said the phone was
plenty fast. ``We've shown it's faster than any competitive phone doing Internet browsing,'' he
said.

There are other, broader technical challenges that affect all carriers as they move to the faster
networks -- both 3G and, coming soon, 4G.

The newer networks are designed to carry bigger chunks of data so that, for example, people can use
their phones to send and receive videos and not just e-mail or text messages. These larger chunks
are not able to travel as far before degrading, however. That presents a serious problem, because
carriers cannot easily erect hundreds of new towers to be closer to users.

Today's cellphone users often just learn to live with the glitches. Rebecca Hwang, 29, a San
Franciscan who has had a 2G iPhone since December 2007, said her calls were cut off periodically and
she did not receive 30 percent of her text messages. But she is still a fan of the phone.

``I love the camera, the music is great, I use the GPS and map all the time,'' Ms.~Hwang said. ``If
I could have a reliable phone, it would be just perfect.''

\pagebreak
\printindex
\end{CJK*}
\end{document}
